\begin{itemize}
\item \textbf{NORMAL.} \emph{Numérica:} Con desaceleración constante v(t) es lineal y x(t) cuadrática: Heun y RK4 reproducen la solución hasta el redondeo (error en d: 8.882$\times10^{-16}$ y 8.882$\times10^{-16}$ m), trapecio y Simpson son exactos y Newton y secante terminan en 2 y 1 iteraciones, frente a 32 de la bisección. Es un caso de verificación, no de comparación de órdenes. \emph{Problema:} La cabina se detiene en 7.5 m, con un margen de 2.5 m respecto a los 10 m disponibles. De ese total, 2.5 m se recorren durante la reacción. La distancia total es más sensible a v$_{0}$: $\kappa$\_rel = 1.673 para v$_{0}$ y 0.3333 para t\_r.
\item \textbf{ESTRÉS.} \emph{Numérica:} Con resistencia cuadrática y h = 0.5 s, Heun comete 0.0819372 m de error en la distancia y RK4 5.946$\times10^{-4}$ m; al refinar el paso los órdenes observados se acercan a 2 y 4. Newton converge cuadráticamente (6 iteraciones), la secante superlinealmente (6) y la bisección linealmente (42). Romberg alcanza la tolerancia con muchas menos evaluaciones que el trapecio. \emph{Problema:} La cabina necesita 21.62 m y excede en 1.62 m la distancia disponible de 20 m. De ese total, 10.5 m se recorren durante la reacción. La distancia total es más sensible a v$_{0}$: $\kappa$\_rel = 1.105 para v$_{0}$ y 0.4856 para t\_r.
\item \textbf{MAL COND..} \emph{Numérica:} La raíz t = T tiene multiplicidad 5: Newton pierde la convergencia cuadrática y avanza con razón 4/5 (104 iteraciones), la secante también se vuelve lineal (146) y el residuo |v(t)| es muchísimo menor que el error en t. Un desplazamiento de 1$\times10^{-12}$ m/s en v mueve la raíz 0.0115416 s (amplificación 1.154$\times10^{10}$): es mal condicionamiento del problema, no fallo de convergencia, porque los tres métodos sí convergen. El t\_det de las EDO tiene un error de 0.109965 s con RK4 aunque la distancia es precisa. \emph{Problema:} La cabina se detiene en 5.833 m, con un margen de 9.17 m respecto a los 15 m disponibles. De ese total, 2.5 m se recorren durante la reacción. La distancia total es más sensible a v$_{0}$: $\kappa$\_rel = 1 para v$_{0}$ y 0.4286 para t\_r. En el frenado plano el instante exacto de detención está mal definido (la cabina «se arrastra» al final), pero la distancia sí es confiable.
\end{itemize}
